% Clase 22 --- La amplitud sin resistencia y su asíntota (§1.2).
% Con R = 0 el denominador 1/C - omega^2 L se anula exactamente en la frecuencia
% natural y la amplitud se va a infinito. La figura muestra las dos ramas y el
% cambio de signo, que es el desfasaje de 180 grados al cruzar la resonancia.
\begin{center}
\begin{tikzpicture}[scale=1]

  \begin{axis}[
    fisfig, width=10.5cm, height=5.4cm,
    xlabel={$\omega$ (frecuencia del generador)}, ylabel={$Q_0$},
    xmin=0, xmax=2.3, ymin=-3.6, ymax=3.6,
    xtick={1}, xticklabels={$\frac{1}{\sqrt{LC}}$},
    ytick={1}, yticklabels={$\varepsilon_0 C$},
  ]
    \addplot[figblue, smooth, samples=200, domain=0:0.93] {1/(1-x^2)};
    \addplot[figblue, smooth, samples=200, domain=1.07:2.3] {1/(1-x^2)};
    \draw[figred, dashed, line width=1pt]
      (axis cs:1,-3.6) -- (axis cs:1,3.6);
    \draw[figaux] (axis cs:0,0) -- (axis cs:2.3,0);
    \node[figlblaux, figred, anchor=south west] at (axis cs:1.04,2.1)
      {resonancia};
    \node[figlblaux, anchor=north east] at (axis cs:2.25,-0.15)
      {$Q_0 \to 0$};
  \end{axis}

  \node[figlbl, anchor=north, align=center] at (5.2,-1.75)
    {$Q_0 = \dfrac{\varepsilon_0}{\dfrac{1}{C} - \omega^2 L}$};

\end{tikzpicture}\\[0.5em]
{\small\itshape Acá $\omega$ es la frecuencia del \textbf{generador} y es
arbitraria: no hay que confundirla con la natural del circuito, que es un dato
de $L$ y $C$. Buscar el régimen permanente es proponer que la carga oscile con
la frecuencia del forzante, $Q = Q_0\cos\omega t$; como cada derivada segunda
trae un $-\omega^2$, la ecuación se vuelve algebraica y sale
$Q_0 = \varepsilon_0/(1/C - \omega^2 L)$ en un renglón. Los dos extremos son
razonables: a frecuencia muy baja el circuito se comporta como si sólo estuviera
el condensador ($Q_0 \to \varepsilon_0 C$, la carga que pondría una batería), y a
frecuencia muy alta la bobina no deja pasar nada y la amplitud se va a cero. Lo
llamativo está en el medio: el denominador se anula exactamente cuando
$\omega = 1/\sqrt{LC}$ y la amplitud \textbf{diverge}. Ésa es la
\textbf{resonancia}, y el cambio de signo al cruzarla no es un error de cuentas:
significa que la carga pasa a oscilar en contrafase con el generador. Que dé
infinito, claro, es una señal de que el modelo sin resistencia ya no sirve
justo ahí.}
\end{center}
